% Section 5: Numerical Experiments
% Word target: ~1500 words

\section{Numerical Experiments}\label{sec:experiments}

This section evaluates the proposed bidirectional meeting point model through
controlled numerical experiments. We compare six model variants across two
scenarios, conduct sensitivity analyses on the key deployment parameters, and
examine service equity across passenger types.

% -----------------------------------------------------------------------
\subsection{Experimental Setup}
% -----------------------------------------------------------------------

\paragraph{Scenarios.}
We use two complementary scenarios. The \emph{synthetic scenario} places
requests uniformly at random on a $20\,\text{km} \times 20\,\text{km}$ grid.
Demand ranges from 100 to 500 requests per horizon; fleet size ranges from 10
to 30 vehicles. Vehicle capacity is 8 passengers and the maximum ride time is
45 minutes. This scenario allows controlled parameter sweeps with full
reproducibility (seeds 42, 43, 44).

The \emph{semi-realistic scenario} is calibrated to a Beijing suburban
district: a $15\,\text{km} \times 15\,\text{km}$ service zone with 80 meeting
points arranged on a $9 \times 9$ grid at 1875\,m spacing, 200 requests
during the morning peak (7--9\,am), and 15 vehicles. This scenario tests
whether the efficiency gains observed in the synthetic setting transfer to a
realistic Chinese city context.

\paragraph{MNL parameters.}
The multinomial logit choice model uses $\beta_{\text{price}} = -0.012$,
$\beta_{\text{time}} = -0.008$, and $\beta_{\text{walk}} = -0.015$, with an
outside option penalty of 5.0. Three passenger types are defined: price-sensitive
(high $|\beta_{\text{price}}|$), time-sensitive (high $|\beta_{\text{time}}|$),
and walk-sensitive (high $|\beta_{\text{walk}}|$), each comprising approximately
one-third of the demand.

\paragraph{Model variants.}
Six variants are evaluated:
\begin{enumerate}
  \item \textbf{DoorToDoor}: no meeting points; vehicles serve origin--destination
        pairs directly.
  \item \textbf{SingleSidedPickup}: meeting points at pickup only; door-to-door
        dropoff.
  \item \textbf{BidirectionalNoChoice}: bidirectional meeting points assigned
        deterministically (no MNL choice model).
  \item \textbf{FullModel}: the proposed model with bidirectional meeting points
        and MNL passenger choice.
  \item \textbf{AblationNoChoice}: bidirectional meeting points with the MNL
        choice model removed (all offers accepted).
  \item \textbf{AblationNoRollingHorizon}: FullModel without rolling horizon
        re-optimization (greedy static assignment).
\end{enumerate}

\paragraph{Performance metrics.}
Nine metrics are recorded: acceptance rate, vehicle-km (vkm), average and
95th-percentile waiting time, average walking distance, average in-vehicle
time (IVT), detour ratio, fairness index (Gini coefficient of per-type
acceptance rates), and CPU time per horizon. All results are averaged across
three seeds.

% -----------------------------------------------------------------------
\subsection{Main Results: Baseline Comparison}
% -----------------------------------------------------------------------

Table~\ref{tab:main-results} reports mean performance across the six variants
for the synthetic scenario with 200 requests and 15 vehicles.

\begin{table}[h]
\centering
\caption{Performance comparison across model variants (synthetic scenario, 200 requests, 15 vehicles, mean $\pm$ std)}
\label{tab:main-results}
\begin{tabular}{lcccccc}
\toprule
Variant & Accept. & vkm & Avg Wait (s) & Avg Walk (m) & Detour & CPU (s) \\
\midrule
DoorToDoor               & 0.610 & 2603.7 & 5003.5 & 0.0    & 3.63 & 12.3 \\
SingleSidedPickup        & 0.440 & 1199.4 & 3847.1 & 1025.3 & 2.42 & 19.1 \\
BidirectionalNoChoice    & 0.533 & 1721.6 & 4626.7 & 1652.6 & 3.60 & 25.7 \\
FullModel                & 0.208 & 628.5  & 450.3  & 95.9   & 1.00 & 1.4  \\
AblationNoChoice         & 0.398 & 923.9  & 309.9  & 78.5   & 0.76 & 15.7 \\
AblationNoRollingHorizon & 0.252 & 753.4  & 3948.8 & 1814.5 & 6.82 & 5.4  \\
\bottomrule
\end{tabular}
\end{table}

\begin{figure}[h]
\centering
[FIGURE 4 PLACEHOLDER: Bar charts comparing acceptance rate and vehicle-km across 6 variants]
\caption{Baseline comparison: acceptance rate and vehicle-km across model variants.}
\label{fig:baseline-comparison}
\end{figure}

\paragraph{Efficiency.}
The FullModel achieves the lowest vehicle-km (628.5) and the lowest detour
ratio (1.00) among all variants. The raw acceptance rate of FullModel (20.8\%)
is lower than DoorToDoor (61.0\%) because the MNL outside option allows
passengers to reject offers that require excessive walking or waiting. However,
the appropriate efficiency metric is vehicle-km \emph{per accepted trip}:
FullModel achieves $628.5 / 0.2082 = 2383.85$\,vkm per unit acceptance,
compared to $2603.7 / 0.6096 = 3662.33$ for DoorToDoor --- a \textbf{34.9\%}
improvement in vehicle utilization. This confirms that bidirectional meeting
point assignment with passenger choice delivers substantially more efficient
routing even after accounting for the lower acceptance rate.

\paragraph{Role of passenger choice.}
Comparing FullModel to AblationNoChoice isolates the contribution of the MNL
choice model. Without choice (AblationNoChoice), acceptance rises to 39.8\%
but vehicle-km increases to 923.9 and the detour ratio rises to 0.76. The
FullModel's lower acceptance reflects passengers self-selecting out of
unattractive offers, which paradoxically improves routing efficiency by
concentrating service on spatially compatible requests.

\paragraph{Role of rolling horizon.}
AblationNoRollingHorizon reveals the cost of static assignment: average
waiting time explodes to 3948.8\,s (versus 450.3\,s for FullModel) and the
detour ratio reaches 6.82. Without periodic re-optimization, early-assigned
vehicles accumulate detours that cannot be corrected, leading to cascading
delays. The 16.6\% increase in vehicle-km further confirms that rolling
horizon re-optimization is essential for operational efficiency.

\paragraph{Single-sided vs.\ bidirectional.}
SingleSidedPickup reduces vehicle-km by 53.9\% relative to DoorToDoor
(1199.4 vs.\ 2603.7) but imposes a 1025.3\,m average walk at dropoff.
BidirectionalNoChoice achieves a higher acceptance rate (53.3\%) but at
higher vehicle-km (1721.6) and longer waits (4626.7\,s) than FullModel,
because without the choice model the optimizer cannot anticipate which
meeting point assignments passengers will accept.

% -----------------------------------------------------------------------
\subsection{Sensitivity Analysis}
% -----------------------------------------------------------------------

\begin{figure}[h]
\centering
[FIGURE 5 PLACEHOLDER: Sensitivity heatmaps --- walking tolerance and fleet size]
\caption{Sensitivity analysis: (a) acceptance rate vs.\ walking tolerance $\rho$; (b) acceptance rate vs.\ fleet size.}
\label{fig:sensitivity}
\end{figure}

\paragraph{Walking tolerance $\rho$.}
We sweep $\rho \in \{200, 400, 600, 800, 1000\}$\,m with all other parameters
fixed at baseline (200 requests, 15 vehicles). FullModel acceptance rate is
\textbf{0.0\%} at $\rho \leq 400$\,m: the MNL utility function penalizes
walking so heavily that no meeting-point offer is accepted when the required
walk exceeds 400\,m. Acceptance rises monotonically, reaching \textbf{9.0\%}
at $\rho = 1000$\,m. DoorToDoor holds steady at approximately 58\% regardless
of $\rho$, since it imposes no walking requirement.

The implication is clear: $\rho = 1000$\,m is the \emph{minimum viable
threshold} for bidirectional meeting point service. Below this threshold,
the system degenerates to near-zero acceptance and the efficiency gains
disappear. Operators must verify that the target population is willing to
walk at least 1000\,m before deploying bidirectional meeting points.

\paragraph{Fleet size.}
We sweep $n_{\text{vehicles}} \in \{5, 10, 15, 20, 25, 30\}$ with
$n_{\text{requests}} = 200$ fixed. FullModel acceptance peaks at
\textbf{24.0\%} with $n_{\text{vehicles}} = 30$ (15 vehicles per 100
requests). The marginal gain in acceptance rate diminishes beyond 20
vehicles: adding vehicles 21--30 yields only 2.1 percentage points of
additional acceptance, compared to 8.3 pp for vehicles 11--20. This
diminishing returns pattern indicates that vehicle scarcity is the binding
constraint at low fleet sizes, while at higher fleet sizes the MNL outside
option becomes the binding constraint.

At $n_{\text{vehicles}} = 15$ (baseline), DoorToDoor acceptance exceeds
FullModel by 38 percentage points (61.0\% vs.\ 20.8\%), confirming that
door-to-door service is more robust under vehicle scarcity because it does
not impose walking penalties. Operators should therefore ensure adequate
fleet supply before deploying bidirectional meeting points.

% -----------------------------------------------------------------------
\subsection{Equity Analysis}
% -----------------------------------------------------------------------

Table~\ref{tab:main-results} reports aggregate acceptance rates, but these
mask heterogeneity across passenger types. We disaggregate FullModel
acceptance by type using the equity analysis from Phase 4.

\paragraph{Per-type acceptance rates.}
Price-sensitive passengers achieve the highest acceptance rate at
\textbf{25.4\%}, followed by walk-sensitive passengers at \textbf{20.2\%},
and time-sensitive passengers at \textbf{14.4\%}. The 11 percentage point
gap between price-sensitive and time-sensitive passengers reflects a
systematic disadvantage: time-sensitive passengers face longer in-vehicle
travel times (avg IVT 1109.6\,s) and higher waiting times (avg 603.2\,s)
relative to other types, making meeting-point offers less attractive under
their utility function.

\paragraph{Gini coefficient.}
The Gini coefficient of per-type acceptance rates is \textbf{0.1216},
indicating moderate inequality. Approximately 12\% of the total acceptance
rate variation is attributable to passenger type rather than random demand
variation. While this level of inequality is lower than that observed in
DoorToDoor (Gini = 0.5385) and BidirectionalNoChoice (Gini = 0.4587), it
remains a policy concern: time-sensitive passengers --- who may include
commuters with rigid schedules or passengers with medical appointments ---
are systematically less well served by the bidirectional meeting point system.

\paragraph{Implications.}
The equity analysis motivates Recommendation R3 in Section~\ref{sec:policy}:
operators should monitor per-type acceptance rates and consider targeted
interventions (fare subsidies, door-to-door fallback) for disadvantaged
passenger segments. The Gini coefficient provides a single scalar metric
for regulatory reporting.
